% Folder 136, batch 09 — archivists' pages 161–180.
% Pass: Opus 5.5 (claude-opus-5-5), 2026-10-03 — first pass, unchecked against the pages by a human.
%
% The odd pages 161–179 and the verso 162 are in his hand. The other even
% pages (164–180) are typed pages of a SGA 7 exposé IX typescript, unrelated
% to the mathematics in hand, and are skipped.
\documentclass[11pt,a4paper]{article}
\input{../preamble/grothendieck}

\folder{136}
\batch{9}
\pages{161}{180}
\foldertitle{Complexe de De Rham à puissance divisée [conférence de 1976 à l’IHÉS] : notes manuscrites (s.d.).}
\dating{[à partir de 1975-1976]}
\watermark{Édition de démonstration}

\begin{document}

\section{Opérations sur complexes de De Rham à p.d.}

\page{161}
\note{Le titre ci-dessus est celui qu'il inscrit, au crayon, dans un cadre en tête de la page. Les notations $\mathcal{A}(S,S^{+},t)_{*}$, $\mathcal{A}_{S}$, $\beta_{S'}$ viennent des pages qui précèdent ce lot.}

\struck{Proposition} \textbf{Théorème.} Soient $(S,S^{+},t)$ et $(S',S'^{+},t')$ deux anneaux (comm.) à idéal à p.d.\ et élément $t$ (\uncertain{OPS} $S^{+}=(t)_{\mathrm{pd}}$, $S'^{+}=(t')_{\mathrm{pd}}$), et $\mathcal{A}(S,S^{+},t)_{*}$ et $\mathcal{A}(S',S'^{+},t')_{*}$ \ill{} qu'on considère comme ann.\ \ill{} ($=\mathcal{A}_{S}$, $\mathcal{A}_{S'}$) \uncertain{affinables} ½simpliciaux (faisant abstraction de leurs structures graduées, filtrées, à p.d., et de \ill{}). Alors les hom.\ de ann.\ \uncertain{affinables} ½simpliciaux $\varphi_{*}$ correspondent biunivoquement aux systèmes $(\varphi_{0}, A'_{1}, (B'_{i})_{i\geq 1})$, où
\[
\begin{cases}
\varphi_{0}\colon S\to S' \ \text{un hom.\ d'anneaux}\\
A'_{1}\in S'\{X_{0}\}^{++}=\operatorname{Ker}\bigl(\alpha_{S'}\colon S'\{X_{0}\}\to S'\bigr)\\
B'_{i}\in S'\{X_{0}\}
\end{cases}
\]
($\alpha(X_{0}^{(i)})=0$ $\forall i\geq 1$), donnés en fonction de $\varphi_{*}$ par
\[
\begin{cases}
\varphi_{0} = \text{composante en degré } 0 \text{ de } \varphi_{*},\quad
\varphi_{0}\colon (\mathcal{A}_{S})_{0}\to(\mathcal{A}_{S'})_{0},\\
A'_{1} = k_{0}\,\varphi_{1}(X_{0}),\\
B'_{i} = k_{1}\,\varphi_{1}(X_{0}^{(i)}),
\end{cases}
\]
avec $(\mathcal{A}_{S})_{0}\simeq S$, $(\mathcal{A}_{S'})_{0}\simeq S'$, $\varphi_{1}\colon S\{X_{0}\}[dX_{0}]\to S'\{X_{0}\}[dX_{0}]$ (${}\simeq(\mathcal{A}_{S})_{1}$, $(\mathcal{A}_{S'})_{1}$), $k_{0}$, $k_{1}$ les projections sur les composantes de degré $0$ resp.\ $1$. Les systèmes $(\varphi_{0}, A'_{1}, B'_{i}\ (i\geq 1))$ \uncertain{satisfaisant} aux conditions
\begin{enumerate}
  \item \struck{\ill{}} $\varphi_{0}(t^{(i)})=\varphi_{0}(t)^{(i)}$, i.e.\ $\varphi_{0}$ commute aux p.d.\ sur $(t)_{\mathrm{pd}}$ (!)
  \item $\beta_{S'}(A'_{1})=\varphi_{0}(t)$, i.e.\ si $A'_{1}=\sum_{j\geq 1}a'_{j}X_{0}^{(j)}$, $\sum_{j\geq 1}a'_{j}t'^{(j)}=\varphi_{0}(t)$
  \item $(2B'_{i})J'=0$ \struck{$(\forall\ldots)$} $\forall i\geq 1$, où $J'$ est l'idéal de $S'$ engendré par les coeff.\ $a'_{j}$ de $A'_{1}$ et les coeff.\ $b'_{jk}$ des $B'_{j}=\sum_{k\geq 0}b'_{jk}X_{0}^{(k)}$.
\end{enumerate}
[NB Si $\varphi_{0}(t)$ est \struck{un diviseur} \add{élément régulier} de $S'$, alors on \uncertain{montre} que $J'$ est engendré par un élément régulier et \ill{} (TSVP)

\marginal{NB dans \uncertain{le th.}, \ill{} l'indice $1$ à $A'_{1}$ \ill{}}

\page{162}
(3) signifie que $2B'_{i}=0$ $\forall i\geq 1$. Il en est de même si $\varphi_{0}(t)$ est \ill{} (ex.\ \ill{}) \ill{} que le coefficient dominant $a'_{1}$ de $A'_{1}(X_{0})=a'_{1}X_{0}+\cdots$ est inversible dans $S'$. Ainsi dans les conditions du \uncertain{cor.~5}, si $(\varphi_{0}(T))_{\mathrm{pd}}=(T')_{\mathrm{pd}}$, i.e.\ $\varphi_{0}(T)=\sum_{i\geq 1}\lambda_{i}T^{(i)}$ avec $\lambda_{1}$ inversible, alors \ill{} (2) \ill{} donne $a'_{1}\equiv 1 \bmod \ill{}$, donc $a'_{1}$ est inversible \ill{} la cond.\ $2B'_{i}=0$ $(i\geq 1)$.]
\note{Suite, sur ce verso, du « TSVP » de la page 161 ; la fin de la remarque est d'une écriture rapide et largement illisible.}

\page{163}
On reconstitue \struck{$\varphi$} $\varphi_{*}$ à partir de $\varphi_{0}$, $A'_{1}$ ($=\sum_{j\geq 1}a'_{j}X_{0}^{(j)}$), $B'_{i}$ $(i\geq 1)$ par
\[
\begin{cases}
\varphi_{1}(\lambda)=\varphi_{0}(\lambda) & \text{si } \lambda\in S\\
\varphi_{1}(X_{0}^{(i)})=A_{1}'^{(i)}+B'_{i}\,dX_{0}\\
\varphi_{1}(dX_{0})=dA'_{1}
\end{cases}
\]
($=A_{1}'^{\,\mathrm{dér}}\,dX_{0}$, où $A_{1}'^{\,\mathrm{dér}}=\sum_{j\geq 1}a'_{j}X_{0}^{(j-1)}$), et $\varphi_{n}$ (pour $n\geq 2$) en posant
\[
U_{0}=X_{0},\ U_{1}=X_{0}+X_{1},\ \ldots,\ U_{n-1}=X_{0}+X_{1}+\cdots+X_{n-1}.
\]
\uncertain{Cela notant} que $(\mathcal{A}_{S})_{n}=S\{U_{0},U_{1},\ldots,U_{n}\}[dU_{0},\ldots,dU_{n}]$, par
\[
\begin{cases}
\varphi_{n}(\lambda)=\varphi_{0}(\lambda) & \text{si } \lambda\in S\\
\varphi_{n}(U_{r}^{(i)})=A_{1}'^{(i)}\{U_{r}\}+B'_{i}\{U_{r}\}\,dU_{r} & (0\leq r\leq n-1,\ i\geq 1)\\
\varphi_{n}(dU_{r})=dA'_{1}\{U_{r}\}=A_{1}'^{\,\mathrm{dér}}\{U_{r}\}\,dU_{r} & (0\leq r\leq n-1)
\end{cases}
\]

[NB la condition $2B'_{i}J'=0$ n'intervient pas pour la construction de $\varphi_{1}$, mais \uncertain{seulement} pour $n\geq 2$, pour assurer que si $0\leq r\neq s\leq n-1$, $\varphi_{n}(U_{r}^{(i)})$ \struck{et $\varphi_{n}(dU_{s})$} \add{commute} aux $\varphi_{n}(U_{s}^{(i)})$ et à $\varphi_{n}(dU_{s})$ \ldots]

\textbf{Corollaire 1.} Si $2$ est inversible dans $S'$ \add{et sous l'une des 3 conditions de la remarque après le th.\ (cf.\ « TSVP »)}, alors \struck{$\varphi_{*}$ respecte} $\varphi_{*}$ respecte la structure graduée ; si de plus $S^{+}=(t)_{\mathrm{pd}}$, $\varphi_{*}$ respecte la structure à p.d.\ \struck{Supposons $S^{+}=(t)_{\mathrm{pd}}$.}

\textbf{Corollaire 2.} (\uncertain{Pour} que $\varphi$ commute aux p.d.,) il f.\ et s.\ qu'on ait
\[
B'_{i}=A_{1}'^{(i-1)}B'_{1}\qquad\forall i\geq 1
\]
(donc les $\varphi_{*}$ en question sont donnés par les systèmes $\varphi_{0}$, $A'_{1}$, $B'_{1}$ satisfaisant 1), 2) et $2B'_{1}J'=0$ (qui tient lieu de 3)).

Pour que $\varphi$ respecte les graduations, il f.\ et s.\ que $B'_{i}=0$ $\forall i\geq 1$ — \struck{(}donc \uncertain{des} des $\varphi$

\page{165}
commutant aux p.d.\ — et ils sont donnés par les couples $(\varphi_{0},A'_{1})$ satisfaisant 1) et 2) [\ill{} \ill{} 3)] \struck{cette condition} \struck{toujours vérifiée si $2$ est inversible dans $S'$}

\textbf{Corollaire 3.} Pour que $\varphi_{*}$ \struck{respecte} respecte la filtration \add{(croiss.\ gross.)}, il f.\ et s.\ que $A'_{1}\in\mathrm{Fil}^{1}$ \struck{et} $B'_{i}\in\mathrm{Fil}^{i-1}$ $\forall i\geq 1$, i.e.
\[
A'_{1}=\underbrace{a'_{1}}_{\lambda}X_{0},\qquad
B'_{i}=b'_{i0}+b'_{i1}X_{0}+\cdots+b'_{i,i-1}X_{0}^{(i-1)}
\]
($a'_{1}, b'_{ij}\in S'$) [donc la condition 2) devient
\[
(2')\quad a'_{1}t'=\varphi_{0}(t)
\]
et (3)
\[
(3')\quad (2b'_{ij})J'=0\qquad (i\geq 1,\ 0\leq j\leq i-1)\ ].
\]

Pour que $\varphi_{*}$ \struck{respecte} respecte les graduations \emph{et} les filtrations, il f.\ et s.\ qu'on ait
\[
A'_{1}=a'_{1}X_{0},\qquad B'_{i}=0\quad i\geq 0,
\]
donc, si de plus on suppose $S^{+}=(t)_{\mathrm{pd}}$, ces $\varphi$ correspondent exactement aux couples $(\varphi_{0},\lambda)\colon (S,S^{+},t)\to(S',S'^{+},t')$ \struck{\ill{}}, par prolongement canonique $\mathcal{A}_{S*}\to\mathcal{A}_{S'*}$ !
\note{« $B'_{i}=0$, $i\geq 0$ » : sic, les $B'_{i}$ n'étant définis que pour $i\geq 1$.}

\textbf{Cor.\ 4.} $\varphi_{*}$ respecte \emph{toujours} la filtration décroissante associée au degré extérieur.

Supposons maintenant
\[
S=k\{T\},\ S^{+}=(T)_{\mathrm{pd}},\ t=T;\qquad
S'=k'\{T\},\ S'^{+}=(T)_{\mathrm{pd}},\ t'=T.
\]

(1) La donnée de $\varphi_{0}$ correspond donc à la donnée de :

(a) $\varphi_{00}\colon k\to S'=k'\{T\}$ \struck{\ill{}}

(\uncertain{explicitable} par un hom.\ $k\xrightarrow{\delta_{0}}k'$ \ill{} et des \add{op.\ diff.} $\delta_{i}$ $(i\geq 1)$, $\delta_{i}\colon k\to k'$, définis par $\varphi_{00}(\lambda)=\sum_{i\geq 0}\delta_{i}(\lambda)T^{(i)}$ [en d'autres termes $\delta_{i}(1)=0$ pour $i\geq 1$, $\delta_{i}(\lambda\lambda')=\sum_{j+k=i}\delta_{j}(\lambda)\delta_{k}(\lambda')$ \ldots])

\page{167}
(b) $\varphi_{0}(T)\in S'^{+}$ \struck{soit} \ill{} fixée par la condition 2°) ; \struck{$\varphi_{0}(T)=\sum_{i\geq 1}\eta'_{i}T^{(i)}$ $(\eta'_{i}\in k',\ i\geq 1)$}

\struck{NB si on suppose qu'on se donne d'abord $\varphi=\varphi_{00}\colon k\to k'$, et qu'on exige $\varphi_{*}$ $\varphi_{00}$-semi-linéaire, i.e.\ $\varphi_{0}$ $\varphi_{00}$-semi-linéaire, la donnée 1) se réduit à (b).}

(2) $A'_{1}\in S'\{X_{0}\}\simeq k'\{X_{0},X_{1}\}$, $A'_{1}=\sum_{i,j\geq 0}\alpha_{ij}X_{0}^{(i)}X_{1}^{(j)}$. La condition $A'_{1}\in S'\{X_{0}\}^{++}$ signifie $\alpha_{ij}=0$ si $i=0$, i.e.
\[
A'_{1}=\sum_{i\geq 1,\ j\geq 0}\alpha_{ij}X_{0}^{(i)}X_{1}^{(j)}.
\]
On a $\beta_{S'}(A'_{1})=\sum_{i\geq 1}\alpha_{i,0}T^{(i)}$, et la condition donne
\[
\varphi_{0}(T)=\sum_{i\geq 1}\alpha_{i,0}T^{(i)}.
\]
\note{Ici et plus bas il note $X_{1}$ la variable que la page appelle aussi $T$ : $S'\{X_{0}\}=k'\{T\}\{X_{0}\}\simeq k'\{X_{0},X_{1}\}$ ; lecture de l'identification incertaine.}

(3) $B'_{i}\in S'\{X_{0}\}\simeq k'\{X_{0},X_{1}\}$ sont liés par les seules conditions $(2B'_{i})J'=0$.

En résumé :

\textbf{Cor.\ 5.} Les hom.\ d'anneaux aff.\ ½simpliciaux
\[
\mathcal{A}_{k*}=\mathcal{A}(k\{T\},k\{T\}^{+},T)\to\mathcal{A}_{k'*}=\mathcal{A}(k'\{T\},k'\{T\}^{+},T)
\]
sont définis par les systèmes d'objets
\[
\begin{cases}
\varphi_{00}\colon k\to k'\{T\}\ \text{hom.\ d'anneaux}\\
A'_{1}\in k'\{T\}\{X_{0}\}^{++}\simeq k'\{X_{0},X_{1}\}\\
B'_{i}\in k'\{T\}\{X_{0}\}\simeq k'\{X_{0},X_{1}\}\quad (i\geq 1)
\end{cases}
\]
satisfaisant les conditions

a) $A'_{1}\in k'\{T\}\{X_{0}\}^{++}$, i.e.\ $A'_{1}=\sum_{i\geq 1,\ j\geq 0}\alpha'_{ij}X_{0}^{(i)}X_{1}^{(j)}$

b) $(2B'_{i})J'=0$ $(\forall i\geq 1)$, où $J'\subset S_{k'}=k'\{T\}$ est l'idéal de $S_{k'}$ engendré par les coeff.\ $a'_{j}$ et $b'_{ij}$ \struck{\ill{}} des
\[
A'_{1}=\sum_{j\geq 1}a'_{j}X_{0}^{(j)}\quad (a'_{j}\in k'\{T\}=S_{k'}),\qquad
B'_{i}=\sum_{j\geq 0}b'_{ij}X_{0}^{(j)}\quad (b'_{ij}\in k'\{T\}=S_{k'}).
\]
Donc b) équivaut aussi à

(b$'$) $(2b'_{ij})J'=0$ pour $i\geq 1$, $j\geq 0$.

\page{169}
De plus

a) $\varphi_{*}$ commute aux p.d.\ ssi $B'_{i}=A_{1}'^{(i-1)}B'_{1}$ $\forall i\geq 1$

b) $\varphi_{*}$ commute aux graduations extérieures ssi $B'_{i}=0$ $\forall i\geq 1$

c) $\varphi_{*}$ commute aux filtrations croissantes grossières (dont les termes sont des sous-$S$-modules) ssi $A'_{1}=\lambda'X_{0}$ $(\lambda'\in k'\{T\})$ et $B'_{i}=\sum_{0\leq j\leq i-1}b'_{ij}X_{0}^{(j)}$ \struck{\ill{}}, $b'_{ij}\in k'\{T\}$ (\uncertain{i.e.\ au minimum})

d) $\varphi_{*}$ commute aux filtrations \struck{dé}croissantes fines [dont les termes sont des sous-$k$-modules \ill{} (ou $S$-modules) \ill{}]
\[
\mathrm{Fil}^{\mathrm{gross}}_{d}\mathcal{A}_{kn}
=\Bigl\{\textstyle\sum F_{i_{0}\ldots i_{r}}\,dX_{i_{0}}\wedge\cdots\wedge dX_{i_{r}}\Bigr\}
\]
\[
F_{i_{0}\ldots i_{r}}\in k\{X_{0},\ldots,X_{n}\} \text{ de « degré » }\leq d-r,
\]
donc c'est la filtration croissante la plus fine pour laquelle \struck{\ill{}} $k\subset\mathrm{Fil}_{0}$, $X_{r}, dX_{r}\in\mathrm{Fil}_{1}$, $\mathrm{Fil}_{i}\cdot\mathrm{Fil}_{j}\subset\mathrm{Fil}_{i+j}$, et $x\in\mathrm{Fil}_{d}\Rightarrow x^{(i)}\in\mathrm{Fil}_{di}$ ; c'est aussi la filtration croissante associée à la graduation par le degré total de
\[
(\mathcal{A}_{k})_{n}\simeq k\{X_{0},\ldots,X_{n}\}[dX_{0},\ldots,dX_{n}]\big/\textstyle\sum dX_{i}
\]
\ill{} avec p.d.\ et tel que $X_{r}$ et $dX_{r}$ de degré $1$ — \ill{} — \ill{} ce que $(\mathcal{A}_{k})_{n}$ devient \ill{} un anneau gradué \ill{} gradué $S$, avec \uncertain{affinables} \ill{}] \struck{il f.\ et s.\ qu'on ait}
\[
\begin{cases}
\varphi_{0}(k)\subset k'\\
A'_{1}=\alpha'_{0}X_{0},\quad \alpha'_{0}\in k'\\
B'_{i}=\sum\limits_{\substack{j,k\geq 0\\ j+k\leq i-1}}\alpha'_{ijk}X_{0}^{(j)}X_{1}^{(k)}\quad (\forall i\geq 1)
\end{cases}
\]
(cette condition \uncertain{est vide} si $2$ \uncertain{inversible} dans $k'$, \uncertain{donc} les $B'_{i}$ \uncertain{sont nuls} !)

\marginal{NB La filtration \uncertain{grossière} \ill{} la filtration \uncertain{fine} \ill{} les filtrations \ill{}}
\note{La note marginale, écrite le long du bord gauche, est presque entièrement illisible.}

\page{171}
auquel cas on a
\[
\varphi_{0}(T)=\alpha'_{0}T.
\]
\struck{De si on suppose que $\varphi_{0}(k)\subset k'$, alors}

L'hom.\ $\varphi_{*}$ \struck{respecte} respecte le degré total ssi on a
\[
\begin{cases}
\varphi_{0}k\subset k'\\
A'_{1}=\alpha'_{0}X_{0}\quad (\alpha'_{0}\in k')\\
B'_{i}=\sum\limits_{\substack{j,k\geq 0\\ j+k=i-1}}\alpha'_{i,j,k}X_{0}^{(j)}X_{1}^{(k)},\quad i\geq 1
\end{cases}
\]
(ces conditions \uncertain{vides} si $2$ inversible dans $S'$)

L'hom.\ $\varphi_{*}$ commute aux deux graduations (totale et extérieure) \add{il en commute à la graduation} ssi \struck{\ill{}} \uncertain{ext.}\ et la filtration croissante fine, i.e.\ ssi on a
\[
\begin{cases}
\varphi_{0}k\subset k'\\
A'_{1}=\alpha'_{0}X_{0}\quad (\alpha'_{0}\in k')\\
B'_{i}=0\quad (i\geq 1)
\end{cases}
\]
Et donc il commute aux p.d.\ — \uncertain{vous voyez un minimum} !). Donc les $\varphi_{*}$ en question correspondent 1--1 aux \struck{triples} \add{couples} $(\varphi_{00}\colon k\to k',\ \alpha'_{0}\in k')$, et sont les hom.
\[
\mathcal{A}_{*}(S_{k},S_{k}^{+},T)\to\mathcal{A}_{*}(S_{k'},S_{k'}^{+},T)
\]
associés à l'hom.
\[
(u_{\varphi_{00},\alpha'_{0}})\colon (S_{k},S_{k}^{+},T)\to(S_{k'},S_{k'}^{+},T)
\]
défini par
\[
u_{\varphi_{00},\alpha'_{0}}\colon k\{T\}\to k'\{T\}
\]
défini par $\varphi_{00}$ sur $k$ et par $T\mapsto\alpha'_{0}T$.

e) Considérons la filtration décroissante sur $\mathcal{A}_{kn}$ et $\mathcal{A}_{k'n}$ associée au degré total (filtration par des sous-$S$ ou sous-$S'$-modules). Pour que $\varphi_{*}$ la respecte, il f.\ et s.\ que l'on ait
\[
B'_{i}=\sum_{\substack{j,k\geq 0\\ j+k\geq i-1}}\alpha'_{i,j,k}X_{0}^{(j)}X_{1}^{(k)}
\]
(conditions \uncertain{toujours vérifiées}, en particulier si $2$ inversible)
\note{La condition sur $\varphi_{0}$ et $A'_{1}$ qui devrait précéder celle-ci sous e) ne se lit pas sur la page : seule la ligne des $B'_{i}$ est écrite.}

\section{Cas des automorphismes de $\widehat{\mathcal{A}}_{k*}$}
\note{Titre encadré par lui en tête de la page 173.}

\page{173}
Nous nous intéressons maintenant aux automorphismes de la $k$-\uncertain{algèbre} \uncertain{affinable} ½simpliciale $\mathcal{A}_{k*}=\mathcal{A}_{*}(k\{T\},k\{T\}^{+},T)$. Mais \ill{} de résultats simples qu'en complétant en
\[
\widehat{\mathcal{A}}_{k*}=\bigl(\widehat{\mathcal{A}}_{k,n}\bigr),
\]
\[
\widehat{\mathcal{A}}_{k,n}\simeq k\{\{X_{0},\ldots,X_{n}\}\}[dX_{0},\ldots,dX_{n}]\big/\textstyle\sum dX_{i}
\]
\[
\simeq k\{\{T\}\}\{\{X_{0},\ldots,X_{n}\}\}[dX_{0},\ldots,dX_{n}]\big/\bigl(\textstyle\sum X_{i}-T,\ \sum dX_{i}\bigr)_{\mathrm{pd}}
\]
\[
\simeq k\{\{T\}\}\{\{X_{0},\ldots,X_{n-1}\}\}[dX_{0},\ldots,dX_{n-1}].
\]
Les résultats sur les hom.\ d'un $\widehat{\mathcal{A}}_{k,*}$ dans un $\widehat{\mathcal{A}}_{k',*}$ sont les \uncertain{mêmes} que tantôt [\ill{}] \uncertain{surtout} en un $\mathcal{A}_{*}(S,S^{+},t)^{\wedge}$ (où on \struck{\ill{}} \ill{} le cas $S\simeq k\{\{T\}\}$), avec
\[
\mathcal{A}_{n}(S,S^{+},t)^{\wedge}\simeq S\{\{X_{0},\ldots,X_{n}\}\}[dX_{0},\ldots,dX_{n}]
\]
mod idéal à p.d.\ engendré par $\sum X_{i}-t$, $\sum dX_{i}$,
\[
\simeq S\{\{X_{0},\ldots,X_{n-1}\}\}[dX_{0},\ldots,dX_{n-1}]
\]
— mais $A'_{1}$ et les $B'_{i}$ $(i\geq 1)$ sont dans $\mathcal{A}_{1}(S,S^{+},t)^{\wedge}\simeq S\{\{X_{0}\}\}$, au lieu de $S\{X_{0}\}$, i.e.\ les polynômes à p.d.\ en $X_{0}$, coeff.\ dans $S$, qui les exprimaient deviennent des séries à p.d.\ avec une

\page{175}
infinité de termes].

Soit $\varphi_{*}$ un \struck{auto} \uncertain{endomorphisme} de la $k$-algèbre diff.\ ½simpliciale $\widehat{\mathcal{A}}_{k*}$, associé à
\[
A,\ \struck{B}\,B_{i}\in\widehat{S}_{k}\{\{X_{0}\}\}\simeq k\{\{T\}\}\{\{X_{0}\}\}\simeq k\{\{X_{0},X_{1}\}\}\qquad (i\geq 1).
\]
Posons
\[
\begin{cases}
A=\sum\limits_{j\geq 1}a_{j}X_{0}^{(j)} & (a_{j}\in\widehat{S}_{k}\simeq k\{\{T\}\},\ j\geq 0)\\
\text{avec } a_{j}=\sum\limits_{k\geq 0}a_{j,k}T^{(k)}, & a_{jk}\in k \text{ pour } j\geq 1,\ k\geq 0.
\end{cases}
\]
On peut aussi écrire
\[
A=\sum_{i\geq 1,\ j\geq 0}\alpha_{ij}X_{0}^{(i)}X_{1}^{(j)},\qquad \alpha_{ij}\in k.
\]
Posons
\[
\varphi_{0}(T)=\sum_{j\geq 1}\lambda_{j}T^{(j)},\qquad \lambda_{j}\in k.
\]
Ceci posé, un calcul immédiat montre
\[
\boxed{a_{1,0}=\alpha_{1,0}=\lambda_{1}}\in k.
\]
Cet élément de $k$ s'appelle le \struck{discriminant} \add{déterminant} \uncertain{indicateur} (noté $\Delta(\varphi_{*})$) de $\varphi_{*}$ (\uncertain{peut aussi se définir si} $\varphi_{*}\colon\widehat{\mathcal{A}}_{k}\to\widehat{\mathcal{A}}_{k'}$ associé à $k\to k'$ donné d'avance, on trouve un élément de $k'$).
\note{Le mot de la ligne, sous « discriminant » biffé et « déterminant » ajouté au-dessus, se lit « indicateur », doublement souligné et peut-être lui aussi barré ; c'est le terme qu'emploie l'énoncé qui suit.}

\textbf{Théorème.} Pour que $\varphi_{*}$ soit un \emph{automorphisme} de $\widehat{\mathcal{A}}_{k*}$, il f.\ et s.\ que son indicateur soit un élément inversible de $k$. En ce cas, $\Delta(\varphi_{*})$ dépend multiplicativement de $\varphi_{*}$, et $\Delta(1)=1$ — donc $\Delta$ est un hom.\ des groupes
\[
\Omega(k)=\operatorname{Aut}_{k\text{-alg.\ diff.\ }\frac{1}{2}\text{simpl.}}\bigl(\widehat{\mathcal{A}}_{k*}\bigr)
\]
dans $k^{*}=\mathbf{G}_{m}(k)$.

\page{177}
Les \struck{automorphismes de} éléments de $\Omega(k)$ correspondent ainsi aux systèmes $(A,(B_{i})_{i\geq 1})$ d'objets de $\widehat{S}_{k}\{\{X_{0}\}\}$,
\[
A=\sum_{j\geq 1}a_{j}X_{0}^{(j)},\qquad
B_{i}=\sum_{j\geq 0}b_{ij}X_{0}^{(j)}\qquad\Bigm|\qquad a_{j}, b_{ij}\in\widehat{S}_{k}=k\{\{T\}\}
\]
tels que

a) $\varepsilon(a_{1})$ ($=\Delta(\varphi_{*})$) $\in k^{*}$ \quad ($\varepsilon\colon\widehat{S}_{k}\to k$ l'augmentation)

b) $2B_{i}=0$ \add{$\forall i\geq 1$}, i.e.\ $2b_{ij}=0$ $\forall i\geq 1$, $j\geq 0$.

On peut aussi l'écrire comme l'ens.\ des systèmes $(a_{jk})_{j\geq 1,k\geq 0}$, $(b_{ijk})_{i\geq 1,j\geq 0,k\geq 0}$ d'éléments de $k$ (qui \struck{sont les coeff.\ des} \struck{développements de} $a_{j}$ et $b_{ij}$ comme séries à p.d.\ en $T$) avec les seules conditions
\[
a_{10}\in k^{*},\qquad 2b_{ijk}=0.
\]
Pour $k$ variable, le \struck{\ill{}} \add{foncteur} $\Omega(k)$ est donc représentable par \add{le schéma affine sur $\mathbf{Z}$}
\[
\mathbf{G}_{m,\mathbf{Z}}\times{}_{2}\mathbf{G}_{a}^{\mathbf{N}^{*}\times\mathbf{N}}\times{}_{2}\mathbf{G}_{a}^{\mathbf{N}^{*}\times\mathbf{N}\times\mathbf{N}}
\]
\note{Lu tel qu'écrit : l'indice $2$ devant le premier facteur $\mathbf{G}_{a}$ (qui correspond aux $a_{jk}$, $j\geq 1$, sans condition hors $a_{10}$) est peut-être un trait parasite.}
\[
\Omega\simeq\operatorname{Spec}\Bigl(\mathbf{Z}\bigl[(A_{jk})_{j\geq 1,k\geq 0},\ (B_{ijk})_{i\geq 1,j\geq 0,k\geq 0}\bigr]\Bigr)
\]
modulo idéal engendré par les $2B_{ijk}$ et localisé du plus p.r.\ à $A_{1,0}$.
\note{Devant le crochet il trace un signe surmonté d'un point, sur le $\mathbf{Z}$, que nous ne savons pas lire.}

\marginal{Ainsi}
$\Omega$ peut être considéré comme un schéma en groupes sur $\mathbf{Z}$, qui est le schéma des automorphismes du foncteur des $\mathbf{Z}$-algèbres vers les $k$-alg.\ diff.\ \uncertain{affinables} ½simpliciales $k\mapsto\widehat{\mathcal{A}}_{k*}$ (ou encore comme nous

\page{179}
disons par abus de langage le schéma en groupes d'automorphismes de l'anneau différentiel ½simplicial $\widehat{\mathcal{A}}_{\mathbf{Z}*}$). On a un hom.\ de schémas en groupes
\[
\Delta\colon\Omega\to\mathbf{G}_{m}
\]
et on a d'ailleurs une section évidente
\[
\delta\colon\mathbf{G}_{m}\to\Omega
\]
($\delta(\lambda)=(A,(B_{i})_{i\geq 1})$ avec $B_{i}=0$ $\forall i\geq 1$,
\[
A=\lambda X_{0}\quad\text{i.e.}\quad \varphi_{1}(X_{0})=\lambda X_{0},\ \text{d'où}\ \varphi_{1}(X_{1})=\lambda X_{1},
\]
\[
\varphi_{0}(T)=\lambda T,\qquad \varphi_{n}(X_{i})=\lambda X_{i}\quad\forall\, 0\leq i\leq n).
\]
Ainsi $\Omega$ est \struck{écrit} un pr.\ semi-direct de $\mathbf{G}_{m}$ et d'un groupe qui est \ill{} unipotent fibre par fibre.

Si $2$ est inversible \add{dans $k$}, \struck{\ill{}} \uncertain{ou plus généralement} \ill{} \uncertain{significations} dans les automorphismes de $\widehat{\mathcal{A}}_{k*}$ respectent le degré extérieur, — donc leurs effets sur $\operatorname{Hom}(X,\widehat{\mathcal{A}}_{k*})$ \add{$X$ ens.\ ½simpl.}\ \ill{} sur la \uncertain{coh.}, des \uncertain{opérations cohomologiques} \uncertain{ne changeant} \ill{} le degré — i.e.\ de degré $0$. Donc \ill{} \uncertain{idées d'invariance} \ill{} que \ill{} (se \uncertain{factorisant} \ill{} $\Delta$) : \uncertain{mult.}\ par $\lambda^{i}$ dans $H^{i}$ $(\lambda\in k^{*})$
\note{L'argument se poursuit au-delà de ce lot : la phrase est interrompue en bas de la page 179.}

\end{document}
